\documentclass{article}
\usepackage{tekvideoexercises}

\begin{document}
\exercisename{Division i polær form}
\streaklength{5}

\tableofcontents
\newpage

\begin{exercise}{Division i polær form 1-1}

Giv i polær form: $\frac{6(\cos(\frac{\pi}{3}) + i\sin(\frac{\pi}{3}))}{2(\cos(\frac{\pi}{6}) + i\sin(\frac{\pi}{6}))}$.

\answerbox{3\left(\cos\left(\frac{\pi}{6}\right) + i\sin\left(\frac{\pi}{6}\right)\right)}

\hint

Når man dividerer to tal i polær form, divideres moduli, og argumenterne trækkes fra hinanden.

\hint

Regel: $\frac{r(\cos \theta + i \sin \theta)}{s(\cos \phi + i \sin \phi)} = \frac{r}{s} \left(\cos(\theta - \phi) + i \sin(\theta - \phi)\right)$

\hint

Her er $r = 6$, $\theta = \frac{\pi}{3}$, $s = 2$, $\phi = \frac{\pi}{6}$:

\hint
\[
 \frac{6}{2} = 3
\]

\hint
\[
 \\frac{\pi}{3} - \frac{\pi}{6} = \frac{\pi}{6}
\]

\end{exercise}

\newpage

\begin{exercise}{Division i polær form 1-2}

Giv i polær form: $\frac{12(\cos(\frac{2\pi}{3}) + i\sin(\frac{2\pi}{3}))}{4(\cos(\frac{\pi}{3}) + i\sin(\frac{\pi}{3}))}$.

\answerbox{3\left(\cos\left(\frac{\pi}{3}\right) + i\sin\left(\frac{\pi}{3}\right)\right)}

\hint

Beregn modulus:
\[
\frac{12}{4} = 3
\]

\hint

Beregn argument:
\[
\frac{2\pi}{3} - \frac{\pi}{3} = \frac{\pi}{3}
\]

\end{exercise}

\newpage

\begin{exercise}{Division i polær form 1-3}

Giv i polær form: $\frac{8(\cos(\frac{5\pi}{4}) + i\sin(\frac{5\pi}{4}))}{2(\cos(\frac{3\pi}{4}) + i\sin(\frac{3\pi}{4}))}$.

\answerbox{4\left(\cos\left(\frac{\pi}{2}\right) + i\sin\left(\frac{\pi}{2}\right)\right)}

\hint

Beregn modulus:
\[
\frac{8}{2} = 4
\]

\hint

Beregn argument:
\[
\frac{5\pi}{4} - \frac{3\pi}{4} = \frac{2\pi}{4} = \frac{\pi}{2}
\]

\end{exercise}

\newpage

\begin{exercise}{Division i polær form 1-4}

Giv i polær form: $\frac{10(\cos(\pi) + i\sin(\pi))}{5(\cos(\frac{\pi}{2}) + i\sin(\frac{\pi}{2}))}$.

\answerbox{2\left(\cos\left(\frac{\pi}{2}\right) + i\sin\left(\frac{\pi}{2}\right)\right)}

\hint

Beregn modulus:
\[
\frac{10}{5} = 2
\]

\hint

Beregn argument:
\[
\pi - \frac{\pi}{2} = \frac{\pi}{2}
\]

\end{exercise}

\newpage

\begin{exercise}{Division i polær form 1-5}

Giv i polær form: $\frac{15(\cos(\frac{7\pi}{6}) + i\sin(\frac{7\pi}{6}))}{3(\cos(\frac{\pi}{6}) + i\sin(\frac{\pi}{6}))}$.

\answerbox{5\left(\cos(\pi) + i\sin(\pi)\right)}

\hint

Beregn modulus:
\[
\frac{15}{3} = 5
\]

\hint

Beregn argument:
\[
\frac{7\pi}{6} - \frac{\pi}{6} = \frac{6\pi}{6} = \pi
\]

\end{exercise}

\newpage

\begin{exercise}{Division i polær form 1-6}

Giv i polær form: $\frac{20(\cos(\frac{3\pi}{2}) + i\sin(\frac{3\pi}{2}))}{4(\cos(\frac{\pi}{2}) + i\sin(\frac{\pi}{2}))}$.

\answerbox{5\left(\cos(\pi) + i\sin(\pi)\right)}

\hint

Beregn modulus:
\[
\frac{20}{4} = 5
\]

\hint

Beregn argument:
\[
\frac{3\pi}{2} - \frac{\pi}{2} = \pi
\]

\end{exercise}

\end{document}
